\documentclass[a4wide]{book}

\usepackage{labs}
\usepackage{listings}

\begin{document}


% ifdef SLIDES

% ifeq ($(firstword $(subst -, , $(SLIDES))),full)
% SLIDES_TRAINING      = $(strip $(subst -slides, , $(subst full-, , $(SLIDES))))
% SLIDES_CHAPTERS      = $($(call UPPERCASE, $(subst  -,_, $(SLIDES_TRAINING)))_SLIDES)
% else
% SLIDES_TRAINING      = $(firstword $(subst -, ,  $(SLIDES)))
% ifeq ($(SLIDES_TRAINING),sysdev)
% SLIDES_TRAINING = embedded-linux
% endif
% ifeq ($(SLIDES_TRAINING),kernel)
% SLIDES_TRAINING = linux-kernel
% endif
% SLIDES_CHAPTERS      = $(filter $(SLIDES)%, $($(call UPPERCASE, $(SLIDES_TRAINING))_SLIDES))
% ifeq ($(SLIDES_CHAPTERS),)
% $(error "No chapter to build, maybe you're building a single chapter whose name doesn't start with a training session name")
% endif
% SLIDES_TYP      = \
% 	$(SLIDES_COMMON_BEFORE) \
% 	$(foreach s,$(SLIDES_CHAPTERS),$(wildcard slides/$(s)/$(s).typ)) \
% 	$(SLIDES_COMMON_AFTER)
% SLIDES_PICTURES = $(call PICTURES,$(foreach s,$(SLIDES_CHAPTERS),slides/$(s))) $(COMMON_PICTURES)
% $(foreach file,$(SLIDES_TYP),$(if $(wildcard $(file)),,$(error Missing file $(file) !)))
% %-slides.pdf: $(VARS) $(SLIDES_TYP) $(SLIDES_PICTURES) $(STYLESHEET2) $(OUTDIR)/last-update.typ
% 	@mkdir -p $(OUTDIR)
% # We generate a .tex file with \input{} directives (instead of just
% # concatenating all files) so that when there is an error, we are
% # pointed at the right original file and the right line in that file.
% 	rm -f $(OUTDIR)/$(basename $@).tex
% 	echo "#include{last-update}" >> $(OUTDIR)/$(basename $@).typ
% 	echo "#include{$(VARS)}" >> $(OUTDIR)/$(basename $@).typ
% 	for f in $(filter %.typ,$^) ; do \
% 		cp $$f $(OUTDIR)/`basename $$f` ; \
% 		sed -i 's%__SESSION_NAME__%$(SLIDES_TRAINING)%' $(OUTDIR)/`basename $$f` ; \
% 		printf "\input{%s}\n" `basename $$f .tex` >> $(OUTDIR)/$(basename $@).typ ; \
% 	done
% 	(cd $(OUTDIR); $(PDFTYPST_ENV) $(PDFTYPST) $(PDFTYPST_OPT) $(basename $@).typ)
% # The second call to pdflatex is to be sure that we have a correct table of
% # content and index
% 	(cd $(OUTDIR); $(PDFTYPST_ENV) $(PDFTYPST) $(PDFTYPST_OPT) $(basename $@).typ > /dev/null 2>&1)
% # We use cat to overwrite the final destination file instead of mv, so
% # that evince notices that the file has changed and automatically
% # reloads it (which doesn't happen if we use mv here). This is called
% # 'Maxime's feature'.
% 	cat out/$@ > $@
% else
% FORCE:
% %-slides.pdf: FORCE
% 	@$(MAKE) $@ SLIDES=$*
% endif